%%%
% src::
%     url = https://tex.stackexchange.com/a/748476
%%%

\documentclass[varwidth, border = 3pt]{standalone}

%%%
% Nous allons utiliser la macro ''\QueryFiles'' du \pack
% ''l3sys-query'' pour interagir avec le \syst de fichiers
% du \os_fr.
%%%
\usepackage{l3sys-query}

\begin{document}

\section{Contenu importé via un motif de recherche}

%%%
% Expliquons les lignes de code qui suivent.
%
%     1) ''sort = name'' demande d'ordonner les fichiers capturés
%     suivant leur nom.
%
%     1) ''pattern'' indique l'usage d'un motif de capture des
%     bons fichiers.
%
%     1) ''content/^[0-9]+.*tex'' est un motif \glob à la saveur
%     \lua. Ici, on recherche dans le dossier path::''content/''
%     les fichiers \ext::''tex'' dont le nom commence par une
%     suite de chiffres.
%
%     1) Le dernier \arg obligatoire contient le code de gestion
%     séquentielle des fichiers capturés. Dans ce code, ''#1''
%     fait \ref au chemin d'un fichier.
%%%
\QueryFiles[sort = name, pattern]{content/^[0-9]+.*tex}{
  \par
  \input{#1}
}

\end{document}
